% section_strong.tex — Section 10: the strong form for cubic graphs — W_4 plus a matching of three-point lines (first solver, 2026-08-19).
% Self-contained; uses \Fp, amsthm environments, Section~\ref{sec:cubic} (thm:cubic, lem:projection, cor:cubic), % lemma_fe.tex — Lemma FE (family equidistribution) for the strong form of the cubic theorem (second solver, 2026-08-19).
% To be \input inside the section on the strong form; uses \Fp, amsthm environments; labels lem:fe.

\begin{lemma}[equidistribution along a family of three-point lines]\label{lem:fe}
Let $p\equiv2\pmod3$, $f(x)=ax^3$ with $a\in\Fp^*$, and let $(i,j)$ be a family of three-point lines with span $J\ge3$ (so that the three residues of a
line are $\alpha_mu$, $m=1,2,3$, with $\alpha_m\in\{-(i+j)/3,(2i-j)/3,(2j-i)/3\}$ nonzero and with pairwise distinct squares $\alpha_m^2$), with direction
$(u^*,v^*)$, $v\equiv ac_1u^3$, $c_1=(i^2-ij+j^2)/3$.  For $\sigma\in\{\pm1\}$ and $m=1,2,3$ let $\Delta_{\sigma,m}(u)=4\sigma/a-3\alpha_m^2u^2$ be the
discriminant of the partner conic $t^2+\alpha_mut+\alpha_m^2u^2-\sigma/a=0$ of the residue $\alpha_mu$, and let $t_{\sigma,m}$ denote one of its roots
(the other is $-\alpha_mu-t_{\sigma,m}$).  Let $F$ be a fixed finite Boolean combination of the following conditions: interval conditions on the least
residues of $u$ (equivalently on the representative $u^*$), of $v(u)$, of the linear forms $\alpha_mu$, and of the roots $t_{\sigma,m}$ and
$-\alpha_mu-t_{\sigma,m}$; and the six membership conditions ``$\Delta_{\sigma,m}(u)$ is a nonzero square''.  Require that $F$ be invariant under each independent root swap $t_{\sigma,m}\leftrightarrow-\alpha_mu-t_{\sigma,m}$, and that root conditions are used only when the corresponding membership holds.  Then
\[
 \#\{u\in\Fp^*:\ F\}\;=\;p\,\mu(F)+O_{i,j}\bigl(\sqrt p\,\log^{16}p\bigr),
\]
where $\mu(F)$ is the probability of $F$ under the following ``local law''.  To specify rational coefficients without suppressing residue branches,
choose a positive common denominator $q$ for the $\alpha_m$, take $T$ uniform on $[0,1)$, and set the normalised least residues of $u$ and $\alpha_mu$
to $\{qT\}$ and $\{q\alpha_mT\}$, respectively.  A chosen representative $u^*$ is obtained by adding the appropriate integer to $\{qT\}$.
The variable $v(u)/p$ is independent uniform, the six membership bits are independent fair coins, and, conditionally on membership,
an averaged choice of each root has independent uniform position $t_{\sigma,m}/p$; its companion has position $\{-q\alpha_mT-t_{\sigma,m}/p\}$.
This law is independent of the chosen common denominator.  The implied constant depends on $(i,j)$ and the fixed Boolean complexity of $F$;
the statement is uniform in interval endpoints, $a$ and the box, and holds for $p$ larger than a constant depending on $(i,j)$.
\end{lemma}
\begin{proof}
First, substituting $x=-(i+j)u/3$ into $(f(x+iu)-f(x))/i$ gives $a(i^2-ij+j^2)u^3/3$, establishing the direction formula.\par
\emph{(1) The fibre products.}  The six discriminants are quadratics in $u$; the roots of $\Delta_{\sigma,m}$ are $\pm(2/\alpha_m)\sqrt{\sigma/(3a)}$.
For fixed $\sigma$ the
three root sets are pairwise disjoint because the $\alpha_m^2$ are pairwise distinct; the root sets of $\Delta_{+,m}$ and $\Delta_{-,m'}$ coincide only if
$\alpha_m^2\equiv-\alpha_{m'}^2\pmod p$, a congruence between fixed rationals which we exclude by taking $p$ large.  Hence the six $\Delta_{\sigma,m}$ are
pairwise coprime squarefree polynomials, their classes in $\overline\Fp(u)^*/(\overline\Fp(u)^*)^2$ are multiplicatively independent, and for every set
$S$ of pairs $(\sigma,m)$ the fibre product $X_S$ of the conics indexed by $S$ over the $u$-line (function field $\overline\Fp(u,\sqrt{\Delta_s}:s\in S)$)
is an absolutely irreducible curve of degree $2^{|S|}$ over the $u$-line, of genus bounded in terms of $|S|$ (at most twelve branch points), whose
$\Fp$-points over a given $u$ number $2^{|S|}$ when all $\Delta_s(u)$, $s\in S$, are nonzero squares and $0$ when one of them is a non-square (the $O(1)$
values of $u$ with some $\Delta_s(u)=0$ are discarded).  Moreover the square roots $\sqrt{\Delta_s}$, $s\in S$, are linearly independent over
$\overline\Fp(u)$, and for $S'$ disjoint from $S$ the function $\prod_{s\in S'}\Delta_s$ is not a square in $\overline\Fp(X_S)$.

\emph{(2) Reduction to character sums.}  Fix a membership pattern, with member set $S$ and non-member set $\bar S$.  Writing
$[\Delta_s(u)\ \text{is a nonzero square}]=\tfrac12(1+\chi(\Delta_s(u)))$ and its negation as $\tfrac12(1-\chi(\Delta_s(u)))$ ($\chi$ the quadratic
character), and every interval condition on a least residue as a Fourier series on $\Z/p$ with $\ell^1$-norm $\le\log p+3$ (or as Selberg polynomials of
degree $H=\lfloor\sqrt p\rfloor$, \cite{Montgomery,Vaaler}), the count of $u$ satisfying $F$ with this membership pattern equals
$2^{-|S|}\sum_{P\in X_S(\Fp)}\prod_{s\in\bar S}\tfrac12(1-\chi(\Delta_s(u(P))))\cdot\prod(\text{interval indicators})$, where the roots $t_s$, $s\in S$, are the
coordinates of $P$.  Expanding, the count is a combination, with total coefficient mass $O(\log^{16}p)$, of the sums
\begin{multline*}
 T(\mathbf h,S')=\sum_{P\in X_S(\Fp)}\chi\Bigl(\prod_{s\in S'}\Delta_s(u)\Bigr)\,
 e\Bigl(\tfrac1p\bigl(h_uu+h_vv(u)+\textstyle\sum_mh_m\alpha_mu\\
 +\sum_{s\in S}(h_st_s+h'_s(-\alpha_{m(s)}u-t_s))\bigr)\Bigr),
 \qquad S'\subset\bar S,\ \|\mathbf h\|_\infty\le H .
\end{multline*}

\emph{(3) Main terms.}  If $S'=\emptyset$ and the additive form is identically zero on $X_S$, then $T=|X_S(\Fp)|=p+O(\sqrt p)$.  The additive form
equals $h_vac_1u^3+(h_u+\sum_mh_m\alpha_m-\sum_sh'_s\alpha_{m(s)})u+\sum_s(h_s-h'_s)t_s$; it vanishes identically iff $h_v\equiv0$ (as $c_1\not\equiv0$),
$h_s=h'_s$ for all $s\in S$ (by the linear independence of the $\sqrt{\Delta_s}$: $\sum_s(h_s-h'_s)t_s=\tfrac12\sum_s(h_s-h'_s)\sqrt{\Delta_s}+$ a polynomial in $u$),
and the coefficient of $u$ is $\equiv0\pmod p$.  The zero-character coefficient for a pattern with $|S|$ members and $6-|S|$ nonmembers is
$2^{-|S|}2^{-(6-|S|)}p=2^{-6}p$: the factor $2^{|S|}$ counts a split fibre, not the points of the irreducible curve.
Thus each of the $2^6$ membership patterns has mass $2^{-6}$, deriving the independent fair bits.
The root averages are normalised Haar measure on the torus with second roots $-\alpha_mu-t_s$.
These frequency vectors are exactly the annihilators of the linear relations of the local law $\mu$
(the residues $\alpha_mu$ are fixed multiples of $u$, and the second root is $-\alpha_mu-t_s$); summing their Fourier coefficients gives $p\,\mu(F)$ up to the
truncation error $O(p/H\cdot\log^{16}p)=O(\sqrt p\log^{16}p)$, and the independence of the memberships and of the roots in $\mu$ is the statement that no other
frequency vector produces a main term.

\emph{(4) Error terms.}  If $S'\ne\emptyset$, the Kummer sheaf of $\chi(\prod_{S'}\Delta_s)$ pulled back to $X_S$ is geometrically non-trivial by (1), and stays
so after tensoring with any Artin--Schreier sheaf (coprime orders); by the Bombieri--Perel'muter bound \cite{Perelmuter,CastroMoreno} $|T(\mathbf h,S')|\le
C(i,j)\sqrt p$ uniformly in $\mathbf h$.  If $S'=\emptyset$ and the additive form is not identically zero on $X_S$, it is a non-constant function of bounded degree
on the absolutely irreducible curve $X_S$ (not of the form $g^p-g+c$ for $p$ large), and Bombieri's bound \cite{Bombieri66} gives $|T|\le C(i,j)\sqrt p$.
Multiplying by the coefficient masses gives the error $O_{i,j}(\sqrt p\log^{16}p)$.
\end{proof}
\begin{remark}
The event that a specified lift lies on no $W_4$ line depends on the unordered pair of partner roots on each slope, so is invariant under every root swap.
Products of these events, which are the only root predicates used in Section~\ref{sec:strong}, have the same invariance.
For intersecting lines, shared residues and coincident conics must use shared membership and root variables.
The pair-count extension is not proved here; it is Hypothesis PC in Section~\ref{sec:strong}.
\end{remark}
 (lem:fe, second solver).
% Data: slack/cube_matching_fast.py (exact counts), slack/g39_agents/{U_model,gamma1,gamma2,gamma2_continuum}.py (constants), logs slack/verification/g39_*.txt.

\section{The strong form for cubic graphs: \texorpdfstring{$W_4$}{W4} plus a matching}\label{sec:strong}

The elementary certificate below is unconditional.  The stronger asymptotic conclusion requires two named hypotheses:
the pair-count law PC and the uniform margin UM.  The quadrature in v1.19 is numerical evidence for them, not a proof.
The unconditional cubic bound remains Theorem~\ref{thm:cubic}.

\subsection*{Three-point lines and the certificate}

We keep the notation of Section~\ref{sec:cubic}: $f(x)=ax^3$, $p\equiv2\pmod3$, box $B=[x_0,x_0+2p)\times[y_0,y_0+2p)$, $P=P(f,B)$, lifts
$(r_x+\delta p,\,s_x+\varepsilon p)$; $W_4$ is the cover ``weight $1$ on every $\pm1$-line with $\ge4$ points'', $L_4$ their number, $S_0\subset P$
the set of points on none of them ($2L_4+|S_0|=\tfrac{11}4p+o(p)$, Theorem~\ref{thm:cubic}).  Three residues $x_1,x_2,x_3$ are collinear on
the residue curve iff $x_1+x_2+x_3=0$ (they are then the roots of $ax^3-mx-c$).  A three-point line of $P$ with a direction of slope $\notin\{0,\infty,\pm1\}$
has three distinct residues; we call it a line of \emph{family $(1,3)$} if its points are $q,\ q+d,\ q+3d$ for an integer vector $d=(u^*,v^*)$: writing
$z=3^{-1}u^*\bmod p$, its residues are $-4z,\,-z,\,5z$ (they sum to $0$) and $v^*\equiv f(-z)-f(-4z)=63az^3\pmod p$; conversely every $u\in\Fp^*$, every lift
of the base residue $-4z$ and every pair of representatives $(u^*,v^*)$ with $q+d,q+3d\in B$ gives such a line, and distinct data give distinct lines.
(In the general description of the three-point lines by steps $(i,j)$ — points $q,q+id,q+jd$, base residue $-(i+j)u/3$ — this is the family $(1,3)$; the
family $(1,2)$, the lines through the four lifts of the centre $(0,f(0))$, is useless for what follows since all its lines meet in four points, and it is
excluded throughout.)  A three-point line is \emph{good}
if its three points lie in $S_0$.

\begin{lemma}[the certificate]\label{lem:cert}
If $M$ is a set of pairwise disjoint good three-point lines, then $\alpha(f,B)\le2L_4+|S_0|-|M|=(\tfrac{11}4+o(1))p-|M|$.
\end{lemma}
\begin{proof}
Weight $1$ on the $\pm1$-lines of $W_4$, weight $1$ on the lines of $M$, weight $1$ on the points of $S_0$ not on a line of $M$: every point of $P$
receives weight $\ge1$, the lines of $M$ are disjoint and inside $S_0$, and a lawful set meets a line in $\le2$ points; the cost is
$2L_4+2|M|+(|S_0|-3|M|)$.
\end{proof}

\begin{lemma}[matching]\label{lem:match}
Let $G$ be the set of good lines of family $(1,3)$ and $\Pi$ the set of unordered pairs of lines of $G$ with a common point.  Then $G$ contains a set of
$|G|-|\Pi|$ pairwise disjoint lines.
\end{lemma}
\begin{proof}
Delete one line from each pair of $\Pi$.
\end{proof}

\subsection*{The local law and the two constants}

Fix a residue $x$ with position $\pi=((x-x_0)\bmod p)/p$ in the $x$-window and thresholds $\theta_+(x)=1-((f(x)-x-y_0+x_0)\bmod p)/p$,
$\theta_-(x)=((f(x)+x-x_0-y_0)\bmod p)/p$ (Section~\ref{sec:cubic}).  Whether a lift $(\delta,\varepsilon)$ of $x$ is in $S_0$ depends, besides
$(\pi,\theta_+,\theta_-,\delta,\varepsilon)$, only on the following data of $x$: whether the partner conics $t^2+xt+x^2\mp1/a=0$ have roots (the
Legendre symbols of $\Delta_\pm=\pm4/a-3x^2$), and the positions of the roots ($t_2$ and $t_3=-x-t_2$ on each slope), through the classes and the line
sizes of Lemma~\ref{lem:lifts}.  Averaging over these data with their limiting law — fair independent coins for the two symbols, and a uniform independent
position of $t_2$ on each slope — defines the \emph{uncovered probability}
\[
 U(\pi,\theta_+,\theta_-,\delta,\varepsilon;\ x_0/p)\in[0,1],
\]
an explicit piecewise polynomial (it factorises as $P_+\cdot P_-$, the two slopes being independent; for $\delta=\varepsilon$ the factor $P_+$ equals
$\tfrac12$ exactly, since the slope-$(+1)$ line through such a lift has $k_++n_+\ge4$ points iff $k_+=3$).
For completeness we define both model constants by finite sums of bounded integrals, including all multiplicities.
Write $\{x\}=x-\lfloor x\rfloor$, $A=(-4,-1,5)$ and $h=(0,1,3)$.
The box-parameter domain is the full torus $(\xi,\eta)\in[0,1)^2$, where $\xi=\{x_0/p\}$ and $\eta=\{y_0/p\}$;
all constants below have both arguments, with no asserted independence of $\eta$.
For parameters $(Z,W)$ set
\[
 \pi_n=\{A_nZ-\xi\},\quad F_n=\{A_n^3W\},\quad
 \theta_{+,n}=1-\{F_n-\pi_n-\eta\},\quad
 \theta_{-,n}=\{F_n+\pi_n-\eta\}.
\]
Given an anchor role $r$, its lifts $d_r,e_r\in\{0,1\}$ and a real direction $(s,v)$, define
\[
 X_n=\xi+\pi_r+d_r+(h_n-h_r)s,\qquad
 Y_n=\eta+\{F_r-\eta\}+e_r+(h_n-h_r)v,
\]
\[
 d_n=X_n-\xi-\pi_n,\qquad e_n=Y_n-\eta-\{F_n-\eta\}.
\]
A line is admissible precisely when all $d_n,e_n\in\{0,1\}$; these differences are integers in the parameterisations below.
Boundary choices have measure zero.  Write $I(L)$ for this indicator and
$U_n(L)=U(\pi_n,\theta_{+,n},\theta_{-,n},d_n,e_n;\xi)$ on admissible data.
The probability $U$ can be evaluated without root labels: for each slope toss a fair membership bit; on membership take
$t\in[0,1)$ uniformly and its companion $\{-A_nZ-t\}$, convert these two absolute root residues to window positions by subtracting $\xi$ modulo one, form the three root classes by the thresholds above,
and use the lift table of Lemma~\ref{lem:lifts} to test whether that lift's line has at most three points.
Average the two-slope indicator.  This is a bounded, explicitly integrable function and is root-swap invariant.

For one line integrate $s\in(-2/3,2/3)$, $w\in[0,1)$, average $\rho\in\{0,1,2\}$, and sum $d_1,e_1\in\{0,1\}$ and $b\in\mathbb Z$.
Use $Z=(s+\rho)/3$, $W=w$, anchor $r=1$, and $v=\{63w\}+b$.  Then
\begin{equation}\label{eq:gamma1}
 \gamma_1(\xi,\eta)=\frac13\sum_{\rho,d_1,e_1,b}\int_{-2/3}^{2/3}\!\int_0^1
 I(L)\prod_{n=1}^3U_n(L)\,dw\,ds .
\end{equation}
Only $|v|<2/3$ contributes, so the integer sum is finite.

For pairs use the following complete list of distinct shared-role configurations and ratios $Z'/Z=P/Q$ in lowest terms:
\[
\begin{array}{c|rrrrrr}
(m,m')&(1,2)&(1,3)&(2,1)&(2,3)&(3,1)&(3,2)\\
P/Q&4&-4/5&1/4&-1/5&-5/4&-5
\end{array}
\]
The role of a point on a family-$(1,3)$ line is intrinsic: its two consecutive gaps have ratio $1:2$.
If the same role is shared, the residue parameter is the same, and two different representatives for either coordinate of the direction
would make the union of the role coordinates span at least $2$; hence both cannot be admissible in a half-open box of side $2$.
Thus a distinct pair has $m\ne m'$.  Two distinct lines share at most one point, so summing the three configurations $m<m'$ counts unordered pairs once.
Equivalently, sum the six ordered configurations and divide by two.

For each $m<m'$ integrate $(s,t)\in(-2/3,2/3)\times[0,1)$ and average $\rho\in\{0,1,2\}$ and $j\in\{0,\dots,Q-1\}$.
Put
\[
 Z=(s+\rho)/3,\quad Z'=(PZ+j)/Q,\quad W=\{Q^3t\},\quad W'=\{P^3t\}.
\]
Construct $L_1$ from $(Z,W)$ with anchor $1$, lifts $(d_1,e_1)\in\{0,1\}^2$, direction $(s,\{63W\}+b_1)$.
Construct $L_2$ from $(Z',W')$ with anchor $m'$, the same physical point as role $m$ of $L_1$ (thus the same lifts),
and direction $(\{3Z'\}+b_2,\{63W'\}+b_3)$.
All $b_i\in\mathbb Z$ are summed; admissibility forces both direction coordinates into $(-2/3,2/3)$ and leaves a finite sum.
With $J(L_1,L_2)$ the joint probability that all their points are uncovered, define
\begin{equation}\label{eq:gamma2}
 \gamma_2(\xi,\eta)=\sum_{m<m'}\frac1{3Q}
 \sum_{\rho,j,d_1,e_1,b_1,b_2,b_3}\int_{-2/3}^{2/3}\!\int_0^1
 I(L_1)I(L_2)J(L_1,L_2)\,dt\,ds .
\end{equation}
Here is an explicit convention for shared variables in $J$.
For each distinct conic indexed by $(\sigma,c)$, with residue coefficient $c$ in
$\{A_n\}\cup\{(P/Q)A_n\}$, use one fair membership bit and one uniform root variable.
Identical coefficients reuse both variables.  If coefficients are opposite, their partner pairs are negatives of one another:
reuse the bit and negate both roots.  More generally conics are identified by equality of $(\sigma,c^2)$,
not by their position in a list; no repeated conic receives a fresh random draw.
The common point is tested once.  In the six configurations above there are five distinct coefficients with distinct squares
apart from the identified common coefficient, so
$J=U_m(L_1)\prod_{n\ne m}U_n(L_1)\prod_{n\ne m'}U_n(L_2)$.
For example the coefficient lists for $m<m'$ are
$\{-16,-4,-1,5,20\}$, $\{-4,-1,5,16/5,4/5\}$ and
$\{-4,-1,5,4/5,1/5\}$.
Additional coincidences modulo $p$ occur only at finitely many primes dividing differences or sums of these fixed rational squares;
they are excluded in the asymptotic count, not assigned independent variables.
The substitution $(W,W')=(\{Q^3t\},\{P^3t\})$ shares the cubic parameter; it is the integral form of averaging its $Q^3$ residue digits.
These definitions, including the finite configuration list, specify a model.  Their applicability to joint arithmetic counts is the following hypothesis.

\paragraph{Hypothesis PC (pair counts).}\label{hyp:pair}
For $p\equiv2\pmod3$ tending to infinity, the number $|\Pi|$ of unordered intersecting pairs of good family-$(1,3)$ lines satisfies
$|\Pi|=\gamma_2(\xi,\eta)p+o(p)$, uniformly in $a\in\Fp^*$ and $(x_0,y_0)$.
The corrected one-family lemma below does not by itself establish PC.

% lemma_fe.tex — Lemma FE (family equidistribution) for the strong form of the cubic theorem (second solver, 2026-08-19).
% To be \input inside the section on the strong form; uses \Fp, amsthm environments; labels lem:fe.

\begin{lemma}[equidistribution along a family of three-point lines]\label{lem:fe}
Let $p\equiv2\pmod3$, $f(x)=ax^3$ with $a\in\Fp^*$, and let $(i,j)$ be a family of three-point lines with span $J\ge3$ (so that the three residues of a
line are $\alpha_mu$, $m=1,2,3$, with $\alpha_m\in\{-(i+j)/3,(2i-j)/3,(2j-i)/3\}$ nonzero and with pairwise distinct squares $\alpha_m^2$), with direction
$(u^*,v^*)$, $v\equiv ac_1u^3$, $c_1=(i^2-ij+j^2)/3$.  For $\sigma\in\{\pm1\}$ and $m=1,2,3$ let $\Delta_{\sigma,m}(u)=4\sigma/a-3\alpha_m^2u^2$ be the
discriminant of the partner conic $t^2+\alpha_mut+\alpha_m^2u^2-\sigma/a=0$ of the residue $\alpha_mu$, and let $t_{\sigma,m}$ denote one of its roots
(the other is $-\alpha_mu-t_{\sigma,m}$).  Let $F$ be a fixed finite Boolean combination of the following conditions: interval conditions on the least
residues of $u$ (equivalently on the representative $u^*$), of $v(u)$, of the linear forms $\alpha_mu$, and of the roots $t_{\sigma,m}$ and
$-\alpha_mu-t_{\sigma,m}$; and the six membership conditions ``$\Delta_{\sigma,m}(u)$ is a nonzero square''.  Require that $F$ be invariant under each independent root swap $t_{\sigma,m}\leftrightarrow-\alpha_mu-t_{\sigma,m}$, and that root conditions are used only when the corresponding membership holds.  Then
\[
 \#\{u\in\Fp^*:\ F\}\;=\;p\,\mu(F)+O_{i,j}\bigl(\sqrt p\,\log^{16}p\bigr),
\]
where $\mu(F)$ is the probability of $F$ under the following ``local law''.  To specify rational coefficients without suppressing residue branches,
choose a positive common denominator $q$ for the $\alpha_m$, take $T$ uniform on $[0,1)$, and set the normalised least residues of $u$ and $\alpha_mu$
to $\{qT\}$ and $\{q\alpha_mT\}$, respectively.  A chosen representative $u^*$ is obtained by adding the appropriate integer to $\{qT\}$.
The variable $v(u)/p$ is independent uniform, the six membership bits are independent fair coins, and, conditionally on membership,
an averaged choice of each root has independent uniform position $t_{\sigma,m}/p$; its companion has position $\{-q\alpha_mT-t_{\sigma,m}/p\}$.
This law is independent of the chosen common denominator.  The implied constant depends on $(i,j)$ and the fixed Boolean complexity of $F$;
the statement is uniform in interval endpoints, $a$ and the box, and holds for $p$ larger than a constant depending on $(i,j)$.
\end{lemma}
\begin{proof}
First, substituting $x=-(i+j)u/3$ into $(f(x+iu)-f(x))/i$ gives $a(i^2-ij+j^2)u^3/3$, establishing the direction formula.\par
\emph{(1) The fibre products.}  The six discriminants are quadratics in $u$; the roots of $\Delta_{\sigma,m}$ are $\pm(2/\alpha_m)\sqrt{\sigma/(3a)}$.
For fixed $\sigma$ the
three root sets are pairwise disjoint because the $\alpha_m^2$ are pairwise distinct; the root sets of $\Delta_{+,m}$ and $\Delta_{-,m'}$ coincide only if
$\alpha_m^2\equiv-\alpha_{m'}^2\pmod p$, a congruence between fixed rationals which we exclude by taking $p$ large.  Hence the six $\Delta_{\sigma,m}$ are
pairwise coprime squarefree polynomials, their classes in $\overline\Fp(u)^*/(\overline\Fp(u)^*)^2$ are multiplicatively independent, and for every set
$S$ of pairs $(\sigma,m)$ the fibre product $X_S$ of the conics indexed by $S$ over the $u$-line (function field $\overline\Fp(u,\sqrt{\Delta_s}:s\in S)$)
is an absolutely irreducible curve of degree $2^{|S|}$ over the $u$-line, of genus bounded in terms of $|S|$ (at most twelve branch points), whose
$\Fp$-points over a given $u$ number $2^{|S|}$ when all $\Delta_s(u)$, $s\in S$, are nonzero squares and $0$ when one of them is a non-square (the $O(1)$
values of $u$ with some $\Delta_s(u)=0$ are discarded).  Moreover the square roots $\sqrt{\Delta_s}$, $s\in S$, are linearly independent over
$\overline\Fp(u)$, and for $S'$ disjoint from $S$ the function $\prod_{s\in S'}\Delta_s$ is not a square in $\overline\Fp(X_S)$.

\emph{(2) Reduction to character sums.}  Fix a membership pattern, with member set $S$ and non-member set $\bar S$.  Writing
$[\Delta_s(u)\ \text{is a nonzero square}]=\tfrac12(1+\chi(\Delta_s(u)))$ and its negation as $\tfrac12(1-\chi(\Delta_s(u)))$ ($\chi$ the quadratic
character), and every interval condition on a least residue as a Fourier series on $\Z/p$ with $\ell^1$-norm $\le\log p+3$ (or as Selberg polynomials of
degree $H=\lfloor\sqrt p\rfloor$, \cite{Montgomery,Vaaler}), the count of $u$ satisfying $F$ with this membership pattern equals
$2^{-|S|}\sum_{P\in X_S(\Fp)}\prod_{s\in\bar S}\tfrac12(1-\chi(\Delta_s(u(P))))\cdot\prod(\text{interval indicators})$, where the roots $t_s$, $s\in S$, are the
coordinates of $P$.  Expanding, the count is a combination, with total coefficient mass $O(\log^{16}p)$, of the sums
\begin{multline*}
 T(\mathbf h,S')=\sum_{P\in X_S(\Fp)}\chi\Bigl(\prod_{s\in S'}\Delta_s(u)\Bigr)\,
 e\Bigl(\tfrac1p\bigl(h_uu+h_vv(u)+\textstyle\sum_mh_m\alpha_mu\\
 +\sum_{s\in S}(h_st_s+h'_s(-\alpha_{m(s)}u-t_s))\bigr)\Bigr),
 \qquad S'\subset\bar S,\ \|\mathbf h\|_\infty\le H .
\end{multline*}

\emph{(3) Main terms.}  If $S'=\emptyset$ and the additive form is identically zero on $X_S$, then $T=|X_S(\Fp)|=p+O(\sqrt p)$.  The additive form
equals $h_vac_1u^3+(h_u+\sum_mh_m\alpha_m-\sum_sh'_s\alpha_{m(s)})u+\sum_s(h_s-h'_s)t_s$; it vanishes identically iff $h_v\equiv0$ (as $c_1\not\equiv0$),
$h_s=h'_s$ for all $s\in S$ (by the linear independence of the $\sqrt{\Delta_s}$: $\sum_s(h_s-h'_s)t_s=\tfrac12\sum_s(h_s-h'_s)\sqrt{\Delta_s}+$ a polynomial in $u$),
and the coefficient of $u$ is $\equiv0\pmod p$.  The zero-character coefficient for a pattern with $|S|$ members and $6-|S|$ nonmembers is
$2^{-|S|}2^{-(6-|S|)}p=2^{-6}p$: the factor $2^{|S|}$ counts a split fibre, not the points of the irreducible curve.
Thus each of the $2^6$ membership patterns has mass $2^{-6}$, deriving the independent fair bits.
The root averages are normalised Haar measure on the torus with second roots $-\alpha_mu-t_s$.
These frequency vectors are exactly the annihilators of the linear relations of the local law $\mu$
(the residues $\alpha_mu$ are fixed multiples of $u$, and the second root is $-\alpha_mu-t_s$); summing their Fourier coefficients gives $p\,\mu(F)$ up to the
truncation error $O(p/H\cdot\log^{16}p)=O(\sqrt p\log^{16}p)$, and the independence of the memberships and of the roots in $\mu$ is the statement that no other
frequency vector produces a main term.

\emph{(4) Error terms.}  If $S'\ne\emptyset$, the Kummer sheaf of $\chi(\prod_{S'}\Delta_s)$ pulled back to $X_S$ is geometrically non-trivial by (1), and stays
so after tensoring with any Artin--Schreier sheaf (coprime orders); by the Bombieri--Perel'muter bound \cite{Perelmuter,CastroMoreno} $|T(\mathbf h,S')|\le
C(i,j)\sqrt p$ uniformly in $\mathbf h$.  If $S'=\emptyset$ and the additive form is not identically zero on $X_S$, it is a non-constant function of bounded degree
on the absolutely irreducible curve $X_S$ (not of the form $g^p-g+c$ for $p$ large), and Bombieri's bound \cite{Bombieri66} gives $|T|\le C(i,j)\sqrt p$.
Multiplying by the coefficient masses gives the error $O_{i,j}(\sqrt p\log^{16}p)$.
\end{proof}
\begin{remark}
The event that a specified lift lies on no $W_4$ line depends on the unordered pair of partner roots on each slope, so is invariant under every root swap.
Products of these events, which are the only root predicates used in Section~\ref{sec:strong}, have the same invariance.
For intersecting lines, shared residues and coincident conics must use shared membership and root variables.
The pair-count extension is not proved here; it is Hypothesis PC in Section~\ref{sec:strong}.
\end{remark}


\begin{proposition}[counts; pair assertion conditional on PC]\label{prop:counts}
Uniformly in $a$ and the box, $|G|=\gamma_1(\xi,\eta)p+o(p)$.
Under Hypothesis PC, also $|\Pi|=\gamma_2(\xi,\eta)p+o(p)$.
\end{proposition}
\begin{proof}
For one line the three residues have coefficients $-4,-1,5$ in $z$ (equivalently divide these by $3$ in $u$).
Their squared coefficients are distinct, and uncovered-lift events are invariant under each root swap.
For each of the finite choices of representatives, lifts and membership bits in \eqref{eq:gamma1}, work in the independent torus coordinates
of Lemma~\ref{lem:fe}: the base parameter, the cubic value and one root per present conic.
After subdivision at the finitely many residue-digit boundaries, admissibility and root-threshold comparisons are affine inequalities in these coordinates.
Their boundary cells in a mesh of side $1/M$ have measure $O(1/M)$, uniformly in $(\xi,\eta)$: there are finitely many hyperplane pieces with fixed normals,
and only their offsets vary.  Deterministic linear identities are eliminated before subdivision.
Approximate from inside and outside by unions of mesh cells, symmetrised under the independent root swaps so that Lemma~\ref{lem:fe} applies.
For fixed $M$ it gives the model count with $o_M(p)$ error, uniformly in $a$ and the box; the difference between the two approximations is
$O(p/M)+o_M(p)$.  Let $p\to\infty$ first and then $M\to\infty$.  Integrating the independent partner data gives the three $U$ factors and proves the stated $o(p)$ error.
The second assertion is precisely PC; no pair independence is inferred by merely reducing the number of conics.
\end{proof}

\paragraph{Hypothesis UM (uniform margin).}\label{hyp:margin}
On the entire parameter domain $[0,1)^2$,
\begin{equation}\label{eq:margin}
 \gamma_1(\xi,\eta)-\gamma_2(\xi,\eta)\ge0.098.
\end{equation}
This is unproved.  In particular a sampled minimum and continuity do not establish it.

\emph{Numerical evidence, quoted from v1.19.}
The archived \texttt{anc/slack/g39\_agents/} computations and their logs report $\gamma_1$ around $0.1287$--$0.1362$
on the sampled grid $\xi,\eta\in\{0,1/4,1/2,3/4\}$ and $\gamma_2$ around $0.028$--$0.031$ at the table's box positions.
The table's $\gamma_2$ column is the v1.19 finite-dump hybrid estimate (\texttt{gamma2.py}); it averages model probabilities on actual intersecting pairs.  It is distinct from the continuum quadrature (\texttt{gamma2\_continuum.py}) for \eqref{eq:gamma2}.  Neither is an enclosure, and neither the endpoints nor $0.098$ are certified extrema.
The $600\times600$ one-line quadrature and the pair quadrature (whose grid sizes are recorded separately in its log) do not cover all box positions.
The table's discrepancy $0.028-0.0212=0.0068$ is to be compared with $0.098-1/12=0.01466\ldots$, about $2.16$ times as large,
not ten times.  No rigorous quadrature error is known here.
The subtractions $2.652=11/4-0.098$ and $2.657=11/4-0.093$ are arithmetic identities.
The first is an extremal upper coefficient under PC and UM; the second would be one under PC and a proved uniform margin $0.093$.
The former v1.19 ``allowance'' $0.005$ proved neither margin and is withdrawn.

\begin{table}[ht]\centering\scriptsize
\begin{tabular}{lccccccc}
\toprule
$p$, box & lines$/p$ & good$/p$ & $\gamma_1$ & pairs$/p$ & $\gamma_2$ & (good$-$pairs)$/p$ & threshold\\
\midrule
$3203$, $(0,0)$        & $1.730$ & $0.1318$ & $0.1287$ & $0.0212$ & $0.028$ & $0.1105$ & $0.0895$\\
$3203$, $(700,-2100)$  & $1.802$ & $0.1339$ & $0.1363$ & $0.0253$ & $0.030$ & $0.1086$ & $0.0720$\\
$12809$, $(0,0)$       & $1.731$ & $0.1313$ & $0.1287$ & $0.0275$ & $0.028$ & $0.1038$ & $0.0847$\\
$12809$, $(-6404,0)$   & $1.801$ & $0.1365$ & $0.1361$ & $0.0315$ & $0.030$ & $0.1050$ & $0.0830$\\
\bottomrule
\end{tabular}
\caption{Family $(1,3)$: exact counts (\texttt{anc/slack/cube\_matching\_fast.py}) versus the uncertified one-line quadrature and pair hybrid estimates (quoted from v1.19); ``threshold'' $=(2L_4+|S_0|-\tfrac83p)/p$ is
the matching size the finite-$p$ certificate needs to reach $\tfrac43N$ (asymptotically $\tfrac1{12}$).}\label{tab:strong}
\end{table}

\begin{theorem}[permutation cubics, strong form; conditional on PC and UM]\label{thm:strong}
Assume Hypotheses PC and UM.  For $p\equiv2\pmod3$ tending to infinity, $a\in\Fp^*$ and every $2p\times2p$ box $B$,
\[
 \alpha(ax^3,B)\le\bigl(11/4-\gamma_1(\xi,\eta)+\gamma_2(\xi,\eta)+o(1)\bigr)p
 \le(2.652+o(1))p=(1.326+o(1))N.
\]
In particular $\alpha(ax^3,B)<(4/3-0.007)N$ for sufficiently large $p$.
\end{theorem}
\begin{proof}
Under PC, Proposition~\ref{prop:counts} and Lemma~\ref{lem:match} give a matching of at least
$(\gamma_1-\gamma_2-o(1))p$ good lines.  Lemma~\ref{lem:cert} gives the first inequality; UM gives the second.
The final strict inequality follows from $4/3-0.007-1.326=1/3000>0$.
\end{proof}

\begin{corollary}[G3, conditional on PC and UM]\label{cor:g3}
Assume Hypotheses PC and UM.  As $p\to\infty$, uniformly over degree-three polynomials $f\in\Fp[x]$ and all $2p\times2p$ boxes,
$\alpha(f,B)\le(4/3+o(1))\,2p$; for permutation cubics the coefficient is at most $1.326$.
\end{corollary}
\begin{proof}
For non-permutation cubics use Lemma~\ref{lem:projection} and \eqref{eq:nonperm}.
Translate permutation cubics to $ax^3$ and use Theorem~\ref{thm:strong} under PC and UM.
\end{proof}

\begin{remark}
The finite certificate of Lemma~\ref{lem:cert} can always be checked point by point.
The historical greedy matching experiments at $p\le797$ and their larger-family variants, archived in
\texttt{anc/slack/verification/cube\_matching\_fast.txt}, are finite computations, not proofs of either hypothesis.
A joint counting proof and a certified enclosure on the full box domain remain necessary to make G3 unconditional.
Together with PC, a proved uniform margin strictly exceeding $1/12$ would suffice for G3 even if the coefficient $2.652$ were lost.
\end{remark}
